% Deutsche Parallelfassung zu why/problem_statement.tex

\chapter{Problemformulierung}

\section{Von der Motivation zur Formalisierung}

Das vorherige Kapitel verfolgte eine Krümmungsänderung durch Geometrie,
Qualifikationen, Evidenz, Regeln und Autorität. Nun bestimmen wir die formalen
Lücken, die eine kohärente Behandlung dieser Beziehungen verhindern.

\section{Zentraler Mangel bestehender Ansätze}

Bestehende Ansätze bieten leistungsfähige geometrische, räumliche,
objektorientierte und numerische Möglichkeiten. Der zentrale Mangel ist nicht
das Fehlen dieser Möglichkeiten, sondern das Fehlen einer
Alignment-spezifischen Darstellung ihrer Beziehungen. Eine übertragene Kurve
begründet keine dauerhafte Objektidentität. Eine gemessene Position gibt nicht
an, ob sie beabsichtigt, physisch realisiert oder beobachtet ist. Eine
erfolgreiche Prüfung von Nebenbedingungen identifiziert weder den Owner der
Regel noch autorisiert sie eine Änderung.

Im geometrischen und numerischen Bereich bleibt eine verwandte Lücke:
Krümmung, Überhöhung, dynamische Größen, Realisierungseffekte und Optimierung
werden häufig als aufeinanderfolgende Aufgaben statt als ausdrücklich
gekoppelte Abhängigkeiten behandelt. Die Thesis behandelt daher sowohl
semantische Kontinuität als auch rechnerische Kopplung. Nur eines der beiden
Probleme zu lösen, ließe die Krümmungsänderung unvollständig verstanden.

Der ursprüngliche Mangel war bereits innerhalb der Übergangstheorie sichtbar.
Benannte Kurven wurden häufig als getrennte Konstruktionen behandelt, obwohl
sich ihre Krümmungsfunktionen anhand gemeinsamer Randwerte und ihres
Ableitungsverhaltens vergleichen ließen. Dadurch waren vier verschiedene
Dinge schwer zu trennen: die Identität einer Funktionsfamilie, ihre
Parametrisierung, die Randbedingungen einer konkreten Aufgabe und die nach
deren Lösung instanziierte Übergangsgeometrie.

\section{Formale Lücke}

Es bezeichne
\[
\gamma : [0,L] \to \R^2
\]
ein ebenes Alignment.

In klassischen Formulierungen ist die Kurve \(\gamma\) das primäre Objekt,
während die Krümmung als
\[
\kappa(s) = \theta'(s)
\]
abgeleitet wird.

Aus dynamischer Sicht kehrt sich die Situation jedoch um:
\begin{itemize}
\item Krümmung bestimmt die Beschleunigung,
\item Krümmungsänderung bestimmt den Ruck,
\item und folglich bestimmt Krümmung das physikalische Verhalten.
\end{itemize}

Dies legt eine grundlegende Diskrepanz offen zwischen:
\begin{itemize}
\item geometrischer Modellierung (Kurve zuerst),
\item und physikalischer Interpretation (Krümmung zuerst).
\end{itemize}

\section{Fehlende Kopplung der Zustandsvariablen}

Es bezeichnen
\[
\kappa(s), \quad \phi(s)
\]
Krümmung und Überhöhung.

In der Praxis sind diese Variablen durch
\[
a_{\mathrm{eff}}(s) = v^2 \kappa(s) - g \sin\phi(s)
\]
stark gekoppelt.

Trotz dieser Kopplung gilt üblicherweise:
\begin{itemize}
\item \(\kappa\) wird geometrisch entworfen,
\item \(\phi\) wird in einem getrennten Prozess entworfen,
\item und ihre Wechselwirkung wird erst anschließend geprüft.
\end{itemize}

Diese Trennung verhindert eine einheitliche Optimierung des Systems.

Betriebliche Kopplung verlangt keine identische funktionale Identität.
Krümmung und Überhöhung können durch getrennte Gesetze
\[
\kappa(s) \quad\text{und}\quad u(s)
\]
mit eigenen Parametern und Randbedingungen beschrieben werden. Die Wahl eines
gemeinsamen Gesetzes, gekoppelter Parameter oder bewusst getrennter Verläufe
ist eine Modellierungs- und Optimierungsentscheidung. ViennaCurve-Beispiele
begründen diesen Freiheitsgrad, indem sie zeigen, wie Krümmungs- und
Überhöhungsentwicklung getrennt geformt werden können, während ihre gemeinsame
betriebliche Folge weiterhin bewertet wird. Diese Darstellung ist eine
ursprüngliche Forschungsinterpretation der Thesis; sie behauptet keinen neuen
einheitlichen Kernel-Begriff.

\section{Lücke zwischen Daten und Realität}

Moderne Eisenbahnprojekte stützen sich stark auf reale Daten, darunter:
\begin{itemize}
\item Messpolylinien,
\item übernommene Alignment-Beschreibungen,
\item und heterogene Koordinatensysteme.
\end{itemize}

Bestehende Modelle haben Schwierigkeiten:
\begin{itemize}
\item verrauschte und unvollständige Daten zu integrieren,
\item unterschiedliche Koordinatenreferenzsysteme in Beziehung zu setzen,
\item und über Maßstäbe hinweg eine konsistente Repräsentation zu erhalten.
\end{itemize}

Dies führt zu einer Trennung zwischen:
\begin{itemize}
\item theoretischen Alignment-Modellen,
\item und realen Ingenieurdaten.
\end{itemize}

\section{Fehlen eines einheitlichen Optimierungs-Frameworks}

Probleme des Alignment-Entwurfs besitzen natürlicherweise die Form
\[
\min J(\kappa,\phi)
\]
unter geometrischen und ingenieurtechnischen Nebenbedingungen.

In der gegenwärtigen Praxis gilt jedoch:
\begin{itemize}
\item Die Zielfunktion ist implizit oder fragmentiert,
\item Nebenbedingungen werden in getrennten Stufen angewendet,
\item und die Optimierung wird nicht als ein einziges Problem formuliert.
\end{itemize}

Dies verhindert den Einsatz systematischer numerischer
Optimierungsverfahren.

\section{Zielformulierung}

Wir wollen den Alignment-Entwurf als einheitliches Problem formulieren:
\[
\text{Finde } (\kappa,\phi)
\]
so, dass:
\begin{itemize}
\item geometrische Konsistenz erfüllt ist,
\item ingenieurtechnische Nebenbedingungen eingehalten sind,
\item und ein globales Zielfunktional minimiert wird.
\end{itemize}

Dies erfordert:
\begin{itemize}
\item eine krümmungsbasierte Repräsentation,
\item ein gekoppeltes Zustandsmodell,
\item und eine konsistente Behandlung von Daten und Nebenbedingungen.
\end{itemize}

Für Übergänge umfasst das Ziel die Kompositionsform
\[
T(s)=H_{\mathrm{in}}(s)+C(s)+H_{\mathrm{out}}(s),
\]
in der eine einlaufende Halbwelle, ein optionaler Klothoidenabschnitt und eine
auslaufende Halbwelle aufeinanderfolgende Teilbereiche belegen und jede
Komponente die Länge null besitzen darf. Dies ist das später in
\cref{sec:berlinish-composition} entwickelte mathematische
Berlinish-Framework; es wird als ursprünglicher Forschungsbeitrag und nicht
als Kernel-Autorität dargestellt.

\section{Problemumfang}

Das in dieser Arbeit behandelte Problem umfasst:
\begin{itemize}
\item die Modellierung von Alignments durch Krümmungsfunktionen,
\item die Einbindung der Überhöhung als Zustandsvariable,
\item die Rekonstruktion der Geometrie aus der Krümmung,
\item die Einbindung von Ingenieurregeln als Nebenbedingungen,
\item und die Formulierung von Optimierungsproblemen.
\end{itemize}

Ausdrücklich ausgeschlossen sind:
\begin{itemize}
\item detaillierte Mehrkörper-Fahrzeugdynamik,
\item die vollständige Modellierung jedes Infrastrukturobjekts,
\item und projektspezifische Bauablaufplanung.
\end{itemize}

\section{Forschungsfragen}

Die in der Einleitung formulierte zentrale Frage wird durch fünf operative
Fragen entfaltet:
\begin{itemize}
\item Welche intrinsische Konstruktion unterscheidet ein Alignment unabhängig
  von seinen Koordinaten und Repräsentationen?
\item Wie kann diese Konstruktion durch ausdrückliche Operatoren als
  geometrischer, dynamischer und betrieblicher Zustand realisiert werden?
\item Wie können beabsichtigter Zustand, physische Realisierung,
  Beobachtungen und angenommenes Ingenieurwissen miteinander verbunden
  bleiben, ohne identisch zu werden?
\item Wie lassen sich anwendbare Regeln und ingenieurtechnische Folgen für die
  Berechnung ausdrücken, während Ownership und Entscheidungsautorität
  ausdrücklich bleiben?
\item Welche Repräsentationen und numerischen Verfahren unterstützen
  Rekonstruktion, Bewertung und Optimierung, ohne das Engineering Object neu
  zu definieren?
\end{itemize}

\section{Arbeitshypothese}

Die Arbeitshypothese lautet, dass eine krümmungsgetriebene konstruktive
Darstellung zusammen mit ausdrücklichen Beziehungen für Zustand,
Realisierung, Beobachtung, Wissen und Bewertung eine kohärente Grundlage für
das Alignment-Ingenieurwesen über Repräsentationen hinweg bietet. Innerhalb
dieser Grundlage können gekoppelte Krümmungs- und Überhöhungsmodelle, hybride
Rekonstruktion, ingenieurtechnische Nebenbedingungen und strukturierte
Optimierung ausgewählte Folgen berechenbar machen, ohne ein numerisches
Ergebnis als Ingenieurentscheidung zu behandeln.

% Deutsche Parallelfassung zu alignment_is_not_a_curve.tex

\section{Alignment}

\subsection{Motivation}

In der klassischen Geometrie wird ein Eisenbahngleis häufig als Kurve
\(\gamma(s)\) im Raum repräsentiert.

Diese Repräsentation reicht für Ingenieurzwecke nicht aus, weil sie die
intrinsischen Größen nicht erfasst, die Konstruktion, Dynamik und Realisierung
bestimmen.

Insbesondere sind Krümmung und Überhöhung keine abgeleiteten Eigenschaften,
sondern primäre Entwurfsvariablen.

\subsection{Definition}

\begin{definition}[Alignment]
Ein \textbf{Alignment} ist ein über der Bogenlänge
\(s \in [0,L]\) definiertes strukturiertes Objekt mit folgenden Komponenten:

\begin{itemize}
\item einer Krümmungsfunktion
\[
\kappa : [0,L] \rightarrow \mathbb{R},
\]

\item einer Überhöhungsfunktion
\[
\phi : [0,L] \rightarrow \mathbb{R},
\]

\item einer Anfangspose
\[
(\gamma_0, \theta_0),
\]
\end{itemize}

zusammen mit den induzierten Zustandsvariablen
\[
\theta(s) = \theta_0 + \int_0^s \kappa(\sigma)\,\dd\sigma,
\]
\[
\gamma(s) = \gamma_0 + \int_0^s
(\cos\theta(\sigma), \sin\theta(\sigma))\,\dd\sigma.
\]
\end{definition}

\subsection{Interpretation}

Ein Alignment wird nicht durch seine Position \(\gamma(s)\), sondern durch
seine Krümmung \(\kappa(s)\) definiert.

Die Kurve \(\gamma(s)\) ist eine durch Integration gewonnene abgeleitete
Größe.

Krümmung ist somit die primäre Entwurfsvariable, während Position eine Folge
ist.

\subsection{Zustandsstruktur}

\begin{definition}[Alignment-Zustand]
Der \textbf{Zustand} eines Alignments an der Position \(s\) ist gegeben durch
\[
(\gamma(s), \theta(s), \kappa(s), \phi(s)).
\]
\end{definition}

Dieser Zustand koppelt Geometrie, Orientierung und Dynamik in einer
einheitlichen Repräsentation.

\subsection{Strukturelle Natur}

\begin{proposition}
Ein Alignment ist ein niedrigdimensionales strukturiertes Objekt, auch wenn
seine Realisierung hochdimensionale Daten umfassen kann.
\end{proposition}

Diese Struktur ermöglicht:
\begin{itemize}
\item effiziente Speicherung,
\item robuste Optimierung,
\item sinnvolle Interpretation.
\end{itemize}

\subsection{Beziehung zur Realisierung}

Das Alignment repräsentiert die beabsichtigte Geometrie.

Das physisch gebaute Gleis ist durch den Realisierungsoperator gegeben:
\[
\gamma_{\mathrm{real}} = \mathcal{R}(\kappa, \phi).
\]

Das Alignment existiert damit im Entwurfsraum, während das Gleis im
physischen Raum existiert.

\subsection{Beziehung zu Daten}

Messdaten liefern üblicherweise abgetastete Positionen oder Polylinien.

Ein Alignment muss aus solchen Daten über inverse Probleme rekonstruiert
werden.

\subsection{Zentrale Erkenntnis}

\begin{proposition}
Ein Alignment ist keine Kurve, sondern ein krümmungsgetriebenes generatives
Modell einer Kurve.
\end{proposition}

\subsection{Zusammenfassung}

\begin{summary}
Das Alignment ist das zentrale Objekt der Eisenbahngeometrie.

Es wird durch Krümmung und Überhöhung definiert, aus denen alle geometrischen
und dynamischen Größen abgeleitet werden.

Diese Sicht verschiebt den Schwerpunkt von der Beschreibung von Kurven zur
Definition generativer Strukturen und bildet die Grundlage des
Alignment-Based Information Modelling.
\end{summary}


\section{Zusammenfassung}

\begin{summary}
Gegenwärtige Verfahren des Alignment-Entwurfs leiden unter der Fragmentierung
zwischen Geometrie, Dynamik und Daten.

Diese Arbeit behandelt die Entwicklung eines einheitlichen,
krümmungsgetriebenen Frameworks, das diese Aspekte in einem einzigen
kohärenten Modell integriert.

Ziel ist die konsistente Modellierung, Analyse und Optimierung von
Eisenbahn-Alignments.
\end{summary}
